% Options for packages loaded elsewhere
% Options for packages loaded elsewhere
\PassOptionsToPackage{unicode}{hyperref}
\PassOptionsToPackage{hyphens}{url}
\PassOptionsToPackage{dvipsnames,svgnames,x11names}{xcolor}
%
\documentclass[
  letterpaper,
  DIV=11,
  numbers=noendperiod]{scrartcl}
\usepackage{xcolor}
\usepackage[margin=1in]{geometry}
\usepackage{amsmath,amssymb}
\setcounter{secnumdepth}{5}
\usepackage{iftex}
\ifPDFTeX
  \usepackage[T1]{fontenc}
  \usepackage[utf8]{inputenc}
  \usepackage{textcomp} % provide euro and other symbols
\else % if luatex or xetex
  \usepackage{unicode-math} % this also loads fontspec
  \defaultfontfeatures{Scale=MatchLowercase}
  \defaultfontfeatures[\rmfamily]{Ligatures=TeX,Scale=1}
\fi
\usepackage{lmodern}
\ifPDFTeX\else
  % xetex/luatex font selection
  \setmonofont[]{DejaVu Sans Mono}
\fi
% Use upquote if available, for straight quotes in verbatim environments
\IfFileExists{upquote.sty}{\usepackage{upquote}}{}
\IfFileExists{microtype.sty}{% use microtype if available
  \usepackage[]{microtype}
  \UseMicrotypeSet[protrusion]{basicmath} % disable protrusion for tt fonts
}{}
\makeatletter
\@ifundefined{KOMAClassName}{% if non-KOMA class
  \IfFileExists{parskip.sty}{%
    \usepackage{parskip}
  }{% else
    \setlength{\parindent}{0pt}
    \setlength{\parskip}{6pt plus 2pt minus 1pt}}
}{% if KOMA class
  \KOMAoptions{parskip=half}}
\makeatother
% Make \paragraph and \subparagraph free-standing
\makeatletter
\ifx\paragraph\undefined\else
  \let\oldparagraph\paragraph
  \renewcommand{\paragraph}{
    \@ifstar
      \xxxParagraphStar
      \xxxParagraphNoStar
  }
  \newcommand{\xxxParagraphStar}[1]{\oldparagraph*{#1}\mbox{}}
  \newcommand{\xxxParagraphNoStar}[1]{\oldparagraph{#1}\mbox{}}
\fi
\ifx\subparagraph\undefined\else
  \let\oldsubparagraph\subparagraph
  \renewcommand{\subparagraph}{
    \@ifstar
      \xxxSubParagraphStar
      \xxxSubParagraphNoStar
  }
  \newcommand{\xxxSubParagraphStar}[1]{\oldsubparagraph*{#1}\mbox{}}
  \newcommand{\xxxSubParagraphNoStar}[1]{\oldsubparagraph{#1}\mbox{}}
\fi
\makeatother

\usepackage{color}
\usepackage{fancyvrb}
\newcommand{\VerbBar}{|}
\newcommand{\VERB}{\Verb[commandchars=\\\{\}]}
\DefineVerbatimEnvironment{Highlighting}{Verbatim}{commandchars=\\\{\}}
% Add ',fontsize=\small' for more characters per line
\usepackage{framed}
\definecolor{shadecolor}{RGB}{241,243,245}
\newenvironment{Shaded}{\begin{snugshade}}{\end{snugshade}}
\newcommand{\AlertTok}[1]{\textcolor[rgb]{0.68,0.00,0.00}{#1}}
\newcommand{\AnnotationTok}[1]{\textcolor[rgb]{0.37,0.37,0.37}{#1}}
\newcommand{\AttributeTok}[1]{\textcolor[rgb]{0.40,0.45,0.13}{#1}}
\newcommand{\BaseNTok}[1]{\textcolor[rgb]{0.68,0.00,0.00}{#1}}
\newcommand{\BuiltInTok}[1]{\textcolor[rgb]{0.00,0.23,0.31}{#1}}
\newcommand{\CharTok}[1]{\textcolor[rgb]{0.13,0.47,0.30}{#1}}
\newcommand{\CommentTok}[1]{\textcolor[rgb]{0.37,0.37,0.37}{#1}}
\newcommand{\CommentVarTok}[1]{\textcolor[rgb]{0.37,0.37,0.37}{\textit{#1}}}
\newcommand{\ConstantTok}[1]{\textcolor[rgb]{0.56,0.35,0.01}{#1}}
\newcommand{\ControlFlowTok}[1]{\textcolor[rgb]{0.00,0.23,0.31}{\textbf{#1}}}
\newcommand{\DataTypeTok}[1]{\textcolor[rgb]{0.68,0.00,0.00}{#1}}
\newcommand{\DecValTok}[1]{\textcolor[rgb]{0.68,0.00,0.00}{#1}}
\newcommand{\DocumentationTok}[1]{\textcolor[rgb]{0.37,0.37,0.37}{\textit{#1}}}
\newcommand{\ErrorTok}[1]{\textcolor[rgb]{0.68,0.00,0.00}{#1}}
\newcommand{\ExtensionTok}[1]{\textcolor[rgb]{0.00,0.23,0.31}{#1}}
\newcommand{\FloatTok}[1]{\textcolor[rgb]{0.68,0.00,0.00}{#1}}
\newcommand{\FunctionTok}[1]{\textcolor[rgb]{0.28,0.35,0.67}{#1}}
\newcommand{\ImportTok}[1]{\textcolor[rgb]{0.00,0.46,0.62}{#1}}
\newcommand{\InformationTok}[1]{\textcolor[rgb]{0.37,0.37,0.37}{#1}}
\newcommand{\KeywordTok}[1]{\textcolor[rgb]{0.00,0.23,0.31}{\textbf{#1}}}
\newcommand{\NormalTok}[1]{\textcolor[rgb]{0.00,0.23,0.31}{#1}}
\newcommand{\OperatorTok}[1]{\textcolor[rgb]{0.37,0.37,0.37}{#1}}
\newcommand{\OtherTok}[1]{\textcolor[rgb]{0.00,0.23,0.31}{#1}}
\newcommand{\PreprocessorTok}[1]{\textcolor[rgb]{0.68,0.00,0.00}{#1}}
\newcommand{\RegionMarkerTok}[1]{\textcolor[rgb]{0.00,0.23,0.31}{#1}}
\newcommand{\SpecialCharTok}[1]{\textcolor[rgb]{0.37,0.37,0.37}{#1}}
\newcommand{\SpecialStringTok}[1]{\textcolor[rgb]{0.13,0.47,0.30}{#1}}
\newcommand{\StringTok}[1]{\textcolor[rgb]{0.13,0.47,0.30}{#1}}
\newcommand{\VariableTok}[1]{\textcolor[rgb]{0.07,0.07,0.07}{#1}}
\newcommand{\VerbatimStringTok}[1]{\textcolor[rgb]{0.13,0.47,0.30}{#1}}
\newcommand{\WarningTok}[1]{\textcolor[rgb]{0.37,0.37,0.37}{\textit{#1}}}

\usepackage{longtable,booktabs,array}
\newcounter{none} % for unnumbered tables
\usepackage{calc} % for calculating minipage widths
% Correct order of tables after \paragraph or \subparagraph
\usepackage{etoolbox}
\makeatletter
\patchcmd\longtable{\par}{\if@noskipsec\mbox{}\fi\par}{}{}
\makeatother
% Allow footnotes in longtable head/foot
\IfFileExists{footnotehyper.sty}{\usepackage{footnotehyper}}{\usepackage{footnote}}
\makesavenoteenv{longtable}
\usepackage{graphicx}
\makeatletter
\newsavebox\pandoc@box
\newcommand*\pandocbounded[1]{% scales image to fit in text height/width
  \sbox\pandoc@box{#1}%
  \Gscale@div\@tempa{\textheight}{\dimexpr\ht\pandoc@box+\dp\pandoc@box\relax}%
  \Gscale@div\@tempb{\linewidth}{\wd\pandoc@box}%
  \ifdim\@tempb\p@<\@tempa\p@\let\@tempa\@tempb\fi% select the smaller of both
  \ifdim\@tempa\p@<\p@\scalebox{\@tempa}{\usebox\pandoc@box}%
  \else\usebox{\pandoc@box}%
  \fi%
}
% Set default figure placement to htbp
\def\fps@figure{htbp}
\makeatother





\setlength{\emergencystretch}{3em} % prevent overfull lines

\providecommand{\tightlist}{%
  \setlength{\itemsep}{0pt}\setlength{\parskip}{0pt}}



 


\KOMAoption{captions}{tableheading}
\makeatletter
\@ifpackageloaded{tcolorbox}{}{\usepackage[skins,breakable]{tcolorbox}}
\@ifpackageloaded{fontawesome5}{}{\usepackage{fontawesome5}}
\definecolor{quarto-callout-color}{HTML}{909090}
\definecolor{quarto-callout-note-color}{HTML}{0758E5}
\definecolor{quarto-callout-important-color}{HTML}{CC1914}
\definecolor{quarto-callout-warning-color}{HTML}{EB9113}
\definecolor{quarto-callout-tip-color}{HTML}{00A047}
\definecolor{quarto-callout-caution-color}{HTML}{FC5300}
\definecolor{quarto-callout-color-frame}{HTML}{acacac}
\definecolor{quarto-callout-note-color-frame}{HTML}{4582ec}
\definecolor{quarto-callout-important-color-frame}{HTML}{d9534f}
\definecolor{quarto-callout-warning-color-frame}{HTML}{f0ad4e}
\definecolor{quarto-callout-tip-color-frame}{HTML}{02b875}
\definecolor{quarto-callout-caution-color-frame}{HTML}{fd7e14}
\makeatother
\makeatletter
\@ifpackageloaded{caption}{}{\usepackage{caption}}
\AtBeginDocument{%
\ifdefined\contentsname
  \renewcommand*\contentsname{Table of contents}
\else
  \newcommand\contentsname{Table of contents}
\fi
\ifdefined\listfigurename
  \renewcommand*\listfigurename{List of Figures}
\else
  \newcommand\listfigurename{List of Figures}
\fi
\ifdefined\listtablename
  \renewcommand*\listtablename{List of Tables}
\else
  \newcommand\listtablename{List of Tables}
\fi
\ifdefined\figurename
  \renewcommand*\figurename{Figure}
\else
  \newcommand\figurename{Figure}
\fi
\ifdefined\tablename
  \renewcommand*\tablename{Table}
\else
  \newcommand\tablename{Table}
\fi
}
\@ifpackageloaded{float}{}{\usepackage{float}}
\floatstyle{ruled}
\@ifundefined{c@chapter}{\newfloat{codelisting}{h}{lop}}{\newfloat{codelisting}{h}{lop}[chapter]}
\floatname{codelisting}{Listing}
\newcommand*\listoflistings{\listof{codelisting}{List of Listings}}
\makeatother
\makeatletter
\makeatother
\makeatletter
\@ifpackageloaded{caption}{}{\usepackage{caption}}
\@ifpackageloaded{subcaption}{}{\usepackage{subcaption}}
\makeatother
\makeatletter
\@ifpackageloaded{tcolorbox}{}{\usepackage[skins,breakable]{tcolorbox}}
\makeatother
\makeatletter
\@ifundefined{shadecolor}{\definecolor{shadecolor}{rgb}{.97, .97, .97}}{}
\makeatother
\makeatletter
\makeatother
\makeatletter
\ifdefined\Shaded\renewenvironment{Shaded}{\begin{tcolorbox}[boxrule=0pt, breakable, enhanced, frame hidden, interior hidden, sharp corners]}{\end{tcolorbox}}\fi
\makeatother
\usepackage{bookmark}
\IfFileExists{xurl.sty}{\usepackage{xurl}}{} % add URL line breaks if available
\urlstyle{same}
\hypersetup{
  pdftitle={Palmer Penguins Part 5: Random Forest versus Linear Models},
  pdfauthor={Ronald `Ryy' G. Thomas},
  colorlinks=true,
  linkcolor={blue},
  filecolor={Maroon},
  citecolor={Blue},
  urlcolor={Blue},
  pdfcreator={LaTeX via pandoc}}


\title{Palmer Penguins Part 5: Random Forest versus Linear Models}
\usepackage{etoolbox}
\makeatletter
\providecommand{\subtitle}[1]{% add subtitle to \maketitle
  \apptocmd{\@title}{\par {\large #1 \par}}{}{}
}
\makeatother
\subtitle{Settling the interpretability-performance tradeoff in
ecological modelling}
\author{Ronald `Ryy' G. Thomas}
\date{2025-01-05}
\begin{document}
\maketitle

\renewcommand*\contentsname{Table of contents}
{
\setcounter{tocdepth}{3}
\tableofcontents
}

\begin{figure}[H]

{\centering \includegraphics[width=0.8\linewidth,height=\textheight,keepaspectratio,alt={Group of penguins standing on rocky Antarctic terrain with the ocean and sky in the background}]{media/images/hero.jpg}

}

\caption{Penguins on rocky Antarctic terrain}

\end{figure}%

\emph{Palmer Penguins series, Part 5 (finale). Photo used under open
licence.}

\begin{tcolorbox}[enhanced jigsaw, arc=.35mm, bottomrule=.15mm, breakable, colback=white, colframe=quarto-callout-note-color-frame, left=2mm, leftrule=.75mm, opacityback=0, rightrule=.15mm, toprule=.15mm]
\begin{minipage}[t]{5.5mm}
\textcolor{quarto-callout-note-color}{\faInfo}
\end{minipage}%
\begin{minipage}[t]{\textwidth - 5.5mm}

\vspace{-3mm}\textbf{Palmer Penguins Data Analysis Series}\vspace{3mm}

This is \textbf{Part 5} (finale) of a 5-part series exploring penguin
morphometrics:

\begin{enumerate}
\def\labelenumi{\arabic{enumi}.}
\tightlist
\item
  \href{../08-palmerpenguinspart1/\%20palmerpenguinspart1/}{Part 1: EDA
  and Simple Regression}
\item
  \href{../09-palmerpenguinspart2/\%20palmerpenguinspart2/}{Part 2:
  Multiple Regression and Species Effects}
\item
  \href{../10-palmerpenguinspart3/\%20palmerpenguinspart3/}{Part 3:
  Cross-Validation and Model Comparison}
\item
  \href{../11-palmerpenguinspart4/\%20palmerpenguinspart4/}{Part 4:
  Model Diagnostics and Interpretation}
\item
  \textbf{Part 5: Random Forest versus Linear Models} (this post)
\end{enumerate}

\end{minipage}%
\end{tcolorbox}

\section{Introduction}\label{introduction}

Whether a random forest meaningfully outperforms a linear model remains
unclear until a careful head-to- head comparison is conducted.
Throughout this series, the species-aware linear model has consistently
delivered an R-squared of approximately 0.86, survived cross-validation
with minimal degradation, and passed every diagnostic check devised. The
natural question for the finale was whether a more flexible algorithm
could do substantially better.

The answer is somewhat unexpected. The random forest improved R-squared
by roughly two percentage points, from about 0.86 to about 0.88, despite
having access to additional predictors and making no assumptions about
linearity. That narrow gap prompts careful reflection on when additional
model complexity is worth the cost in interpretability, reproducibility,
and communication.

This post presents the full comparison, examines feature importance from
both perspectives, and concludes the series with practical guidance for
choosing between linear and ensemble methods in ecological research.

\subsection{Motivations}\label{motivations}

\begin{itemize}
\tightlist
\item
  To settle the question of whether machine learning methods
  meaningfully outperform linear regression on this ecological dataset.
\item
  To establish a principled framework for evaluating the
  interpretability-performance tradeoff, not just an intuitive
  judgement.
\item
  To determine whether the random forest identifies the same important
  predictors as the linear model or reveals hidden structure.
\item
  To provide a concrete example of a case where the simpler model is the
  better choice, as a counterpoint to the assumption that more complex
  methods are always superior.
\item
  To summarise the entire five-part journey into a set of practical
  recommendations.
\item
  To produce a reusable template for model comparison applicable to
  future datasets.
\end{itemize}

\subsection{Objectives}\label{objectives}

\begin{enumerate}
\def\labelenumi{\arabic{enumi}.}
\tightlist
\item
  Train a random forest with all available predictors and compare its
  cross-validated RMSE and R-squared against the species-aware linear
  model under identical evaluation conditions.
\item
  Extract and compare feature importance rankings from both models, and
  assess whether they align with biological expectations.
\item
  Develop a decision framework that maps project requirements
  (interpretability, sample size, regulatory context) to an appropriate
  model class.
\item
  Summarise the key findings and methodological lessons from all five
  parts of the series.
\end{enumerate}

Errors and better approaches are welcome; see the Feedback section at
the end.

\begin{center}

\includegraphics[width=0.3\linewidth,height=\textheight,keepaspectratio]{media/images/r-logo.png}
\end{center}

\emph{Visual interlude before the technical content.}

\section{Prerequisites and Setup}\label{prerequisites-and-setup}

This analysis loads the same core packages used throughout the series,
with \texttt{randomForest} and \texttt{caret} providing the ensemble
learning and cross-validation infrastructure.

\begin{Shaded}
\begin{Highlighting}[]
\FunctionTok{library}\NormalTok{(palmerpenguins)}
\FunctionTok{library}\NormalTok{(tidyverse)}
\FunctionTok{library}\NormalTok{(broom)}
\FunctionTok{library}\NormalTok{(caret)}
\FunctionTok{library}\NormalTok{(randomForest)}
\FunctionTok{library}\NormalTok{(knitr)}
\FunctionTok{library}\NormalTok{(patchwork)}

\FunctionTok{theme\_set}\NormalTok{(}\FunctionTok{theme\_minimal}\NormalTok{(}\AttributeTok{base\_size =} \DecValTok{12}\NormalTok{))}
\NormalTok{penguin\_colors }\OtherTok{\textless{}{-}} \FunctionTok{c}\NormalTok{(}
  \StringTok{"Adelie"} \OtherTok{=} \StringTok{"\#FF6B6B"}\NormalTok{,}
  \StringTok{"Chinstrap"} \OtherTok{=} \StringTok{"\#4ECDC4"}\NormalTok{,}
  \StringTok{"Gentoo"} \OtherTok{=} \StringTok{"\#45B7D1"}
\NormalTok{)}

\FunctionTok{data}\NormalTok{(penguins)}
\NormalTok{penguins\_clean }\OtherTok{\textless{}{-}}\NormalTok{ penguins }\SpecialCharTok{|\textgreater{}} \FunctionTok{drop\_na}\NormalTok{()}

\FunctionTok{set.seed}\NormalTok{(}\DecValTok{42}\NormalTok{)}
\end{Highlighting}
\end{Shaded}

We use the same random seed and cross-validation configuration as Part 3
to ensure that the results are directly comparable.

\section{What is the Interpretability-Performance
Tradeoff?}\label{what-is-the-interpretability-performance-tradeoff}

The interpretability-performance tradeoff is the observation that, in
statistical modelling, gains in predictive accuracy often come at the
expense of the analyst's ability to explain how the model arrives at its
predictions. A linear regression produces a simple equation: body mass
equals a weighted sum of morphometric measurements plus species offsets.
Every coefficient has a direct biological interpretation, and confidence
intervals quantify uncertainty in each estimate.

A random forest, by contrast, aggregates predictions from hundreds of
decision trees, each built on a bootstrap sample with a random subset of
predictors. The resulting predictions are often more accurate, but there
is no single equation to inspect. Feature importance scores indicate
which variables matter most, but they do not tell you the direction or
magnitude of each effect in the way that regression coefficients do.

The tradeoff is not purely technical. In regulatory settings, grant
applications, and peer-reviewed publications, reviewers and stakeholders
often require models whose logic can be stated explicitly. A
two-percentage-point improvement in R-squared may not justify the loss
of that transparency.

\section{Getting Started}\label{getting-started}

\subsection{Cross-Validation
Configuration}\label{cross-validation-configuration}

We establish a single cross-validation scheme that both models will
share, eliminating any advantage from different fold assignments.

\begin{Shaded}
\begin{Highlighting}[]
\NormalTok{train\_control }\OtherTok{\textless{}{-}} \FunctionTok{trainControl}\NormalTok{(}
  \AttributeTok{method =} \StringTok{"cv"}\NormalTok{,}
  \AttributeTok{number =} \DecValTok{10}\NormalTok{,}
  \AttributeTok{savePredictions =} \StringTok{"final"}\NormalTok{,}
  \AttributeTok{verboseIter =} \ConstantTok{FALSE}
\NormalTok{)}
\end{Highlighting}
\end{Shaded}

\subsection{Training the Linear
Champion}\label{training-the-linear-champion}

The species-aware linear model uses the same formula that has been our
workhorse since Part 2.

\begin{Shaded}
\begin{Highlighting}[]
\NormalTok{cv\_linear }\OtherTok{\textless{}{-}} \FunctionTok{train}\NormalTok{(}
\NormalTok{  body\_mass\_g }\SpecialCharTok{\textasciitilde{}}\NormalTok{ bill\_length\_mm }\SpecialCharTok{+}\NormalTok{ bill\_depth\_mm }\SpecialCharTok{+}
\NormalTok{    flipper\_length\_mm }\SpecialCharTok{+}\NormalTok{ species,}
  \AttributeTok{data =}\NormalTok{ penguins\_clean,}
  \AttributeTok{method =} \StringTok{"lm"}\NormalTok{,}
  \AttributeTok{trControl =}\NormalTok{ train\_control}
\NormalTok{)}
\end{Highlighting}
\end{Shaded}

\subsection{Training the Random Forest
Challenger}\label{training-the-random-forest-challenger}

The random forest receives all available predictors, including sex and
island, which were not used in the linear model. This gives the ensemble
every opportunity to outperform.

\begin{Shaded}
\begin{Highlighting}[]
\FunctionTok{set.seed}\NormalTok{(}\DecValTok{123}\NormalTok{)}

\NormalTok{cv\_rf }\OtherTok{\textless{}{-}} \FunctionTok{train}\NormalTok{(}
\NormalTok{  body\_mass\_g }\SpecialCharTok{\textasciitilde{}}\NormalTok{ bill\_length\_mm }\SpecialCharTok{+}\NormalTok{ bill\_depth\_mm }\SpecialCharTok{+}
\NormalTok{    flipper\_length\_mm }\SpecialCharTok{+}\NormalTok{ species }\SpecialCharTok{+}\NormalTok{ sex }\SpecialCharTok{+}\NormalTok{ island,}
  \AttributeTok{data =}\NormalTok{ penguins\_clean,}
  \AttributeTok{method =} \StringTok{"rf"}\NormalTok{,}
  \AttributeTok{trControl =}\NormalTok{ train\_control,}
  \AttributeTok{ntree =} \DecValTok{500}\NormalTok{,}
  \AttributeTok{importance =} \ConstantTok{TRUE}
\NormalTok{)}
\end{Highlighting}
\end{Shaded}

The \texttt{importance\ =\ TRUE} flag instructs the algorithm to compute
permutation-based variable importance, which we examine in the feature
importance section below.

\begin{figure}[H]

{\centering \pandocbounded{
\includegraphics[keepaspectratio,alt={Group of penguins near the Antarctic shoreline with water and rocks in the background}]{media/images/ambiance2.jpg}}

}

\caption{Penguins gathered near the shore}

\end{figure}%

\emph{Midpoint pause before the comparison results.}

\section{Head-to-Head Comparison}\label{head-to-head-comparison}

\subsection{Cross-Validated
Performance}\label{cross-validated-performance}

\begin{Shaded}
\begin{Highlighting}[]
\NormalTok{comparison\_table }\OtherTok{\textless{}{-}} \FunctionTok{tibble}\NormalTok{(}
  \AttributeTok{Model =} \FunctionTok{c}\NormalTok{(}\StringTok{"Linear (species)"}\NormalTok{, }\StringTok{"Random Forest"}\NormalTok{),}
  \AttributeTok{RMSE =} \FunctionTok{c}\NormalTok{(}
\NormalTok{    cv\_linear}\SpecialCharTok{$}\NormalTok{results}\SpecialCharTok{$}\NormalTok{RMSE,}
    \FunctionTok{min}\NormalTok{(cv\_rf}\SpecialCharTok{$}\NormalTok{results}\SpecialCharTok{$}\NormalTok{RMSE)}
\NormalTok{  ),}
  \AttributeTok{R\_squared =} \FunctionTok{c}\NormalTok{(}
\NormalTok{    cv\_linear}\SpecialCharTok{$}\NormalTok{results}\SpecialCharTok{$}\NormalTok{Rsquared,}
    \FunctionTok{max}\NormalTok{(cv\_rf}\SpecialCharTok{$}\NormalTok{results}\SpecialCharTok{$}\NormalTok{Rsquared)}
\NormalTok{  ),}
  \AttributeTok{MAE =} \FunctionTok{c}\NormalTok{(}
\NormalTok{    cv\_linear}\SpecialCharTok{$}\NormalTok{results}\SpecialCharTok{$}\NormalTok{MAE,}
    \FunctionTok{min}\NormalTok{(cv\_rf}\SpecialCharTok{$}\NormalTok{results}\SpecialCharTok{$}\NormalTok{MAE)}
\NormalTok{  )}
\NormalTok{)}

\FunctionTok{kable}\NormalTok{(}
\NormalTok{  comparison\_table,}
  \AttributeTok{digits =} \DecValTok{3}\NormalTok{,}
  \AttributeTok{col.names =} \FunctionTok{c}\NormalTok{(}
    \StringTok{"Model"}\NormalTok{, }\StringTok{"RMSE (g)"}\NormalTok{, }\StringTok{"R{-}squared"}\NormalTok{, }\StringTok{"MAE (g)"}
\NormalTok{  ),}
  \AttributeTok{caption =} \StringTok{"Cross{-}validated performance metrics}
\StringTok{    for the linear model and random forest."}
\NormalTok{)}
\end{Highlighting}
\end{Shaded}

The random forest achieves a modestly higher R-squared and lower RMSE,
but the differences are small. In absolute terms, the RMSE improvement
amounts to a reduction of roughly 20-30 grams in prediction error, a
quantity that is likely within the range of field measurement
uncertainty.

\subsection{Fold-by-Fold Variability}\label{fold-by-fold-variability}

\begin{Shaded}
\begin{Highlighting}[]
\NormalTok{fold\_results }\OtherTok{\textless{}{-}} \FunctionTok{bind\_rows}\NormalTok{(}
\NormalTok{  cv\_linear}\SpecialCharTok{$}\NormalTok{resample }\SpecialCharTok{|\textgreater{}}
    \FunctionTok{mutate}\NormalTok{(}\AttributeTok{Model =} \StringTok{"Linear"}\NormalTok{),}
\NormalTok{  cv\_rf}\SpecialCharTok{$}\NormalTok{resample }\SpecialCharTok{|\textgreater{}}
    \FunctionTok{mutate}\NormalTok{(}\AttributeTok{Model =} \StringTok{"Random Forest"}\NormalTok{)}
\NormalTok{)}

\FunctionTok{ggplot}\NormalTok{(}
\NormalTok{  fold\_results,}
  \FunctionTok{aes}\NormalTok{(}\AttributeTok{x =}\NormalTok{ Model, }\AttributeTok{y =}\NormalTok{ Rsquared, }\AttributeTok{fill =}\NormalTok{ Model)}
\NormalTok{) }\SpecialCharTok{+}
  \FunctionTok{geom\_boxplot}\NormalTok{(}\AttributeTok{alpha =} \FloatTok{0.7}\NormalTok{) }\SpecialCharTok{+}
  \FunctionTok{scale\_fill\_manual}\NormalTok{(}
    \AttributeTok{values =} \FunctionTok{c}\NormalTok{(}
      \StringTok{"Linear"} \OtherTok{=} \StringTok{"\#FF6B6B"}\NormalTok{,}
      \StringTok{"Random Forest"} \OtherTok{=} \StringTok{"\#45B7D1"}
\NormalTok{    )}
\NormalTok{  ) }\SpecialCharTok{+}
  \FunctionTok{labs}\NormalTok{(}
    \AttributeTok{title =} \StringTok{"R{-}squared Distribution Across Folds"}\NormalTok{,}
    \AttributeTok{x =} \ConstantTok{NULL}\NormalTok{,}
    \AttributeTok{y =} \StringTok{"R{-}squared"}
\NormalTok{  ) }\SpecialCharTok{+}
  \FunctionTok{theme\_minimal}\NormalTok{() }\SpecialCharTok{+}
  \FunctionTok{theme}\NormalTok{(}\AttributeTok{legend.position =} \StringTok{"none"}\NormalTok{)}
\end{Highlighting}
\end{Shaded}

The boxplots show substantial overlap between the two models' per-fold
R-squared distributions, confirming that the performance difference is
not statistically compelling for this sample size.

\section{Feature Importance Analysis}\label{feature-importance-analysis}

\subsection{Random Forest Importance}\label{random-forest-importance}

\begin{Shaded}
\begin{Highlighting}[]
\NormalTok{rf\_imp }\OtherTok{\textless{}{-}} \FunctionTok{varImp}\NormalTok{(cv\_rf)}\SpecialCharTok{$}\NormalTok{importance }\SpecialCharTok{|\textgreater{}}
  \FunctionTok{rownames\_to\_column}\NormalTok{(}\StringTok{"Variable"}\NormalTok{) }\SpecialCharTok{|\textgreater{}}
  \FunctionTok{arrange}\NormalTok{(}\FunctionTok{desc}\NormalTok{(Overall)) }\SpecialCharTok{|\textgreater{}}
  \FunctionTok{head}\NormalTok{(}\DecValTok{6}\NormalTok{)}

\FunctionTok{kable}\NormalTok{(}
\NormalTok{  rf\_imp,}
  \AttributeTok{digits =} \DecValTok{1}\NormalTok{,}
  \AttributeTok{col.names =} \FunctionTok{c}\NormalTok{(}\StringTok{"Variable"}\NormalTok{, }\StringTok{"Importance"}\NormalTok{),}
  \AttributeTok{caption =} \StringTok{"Variable importance scores from the}
\StringTok{    random forest (permutation{-}based)."}
\NormalTok{)}
\end{Highlighting}
\end{Shaded}

\subsection{Linear Model Coefficients}\label{linear-model-coefficients}

\begin{Shaded}
\begin{Highlighting}[]
\NormalTok{linear\_coefs }\OtherTok{\textless{}{-}} \FunctionTok{tidy}\NormalTok{(}
\NormalTok{  cv\_linear}\SpecialCharTok{$}\NormalTok{finalModel, }\AttributeTok{conf.int =} \ConstantTok{TRUE}
\NormalTok{) }\SpecialCharTok{|\textgreater{}}
  \FunctionTok{filter}\NormalTok{(term }\SpecialCharTok{!=} \StringTok{"(Intercept)"}\NormalTok{) }\SpecialCharTok{|\textgreater{}}
  \FunctionTok{mutate}\NormalTok{(}
    \AttributeTok{abs\_effect =} \FunctionTok{abs}\NormalTok{(estimate)}
\NormalTok{  ) }\SpecialCharTok{|\textgreater{}}
  \FunctionTok{arrange}\NormalTok{(}\FunctionTok{desc}\NormalTok{(abs\_effect)) }\SpecialCharTok{|\textgreater{}}
  \FunctionTok{select}\NormalTok{(term, estimate, conf.low, conf.high)}

\FunctionTok{kable}\NormalTok{(}
\NormalTok{  linear\_coefs,}
  \AttributeTok{digits =} \DecValTok{1}\NormalTok{,}
  \AttributeTok{col.names =} \FunctionTok{c}\NormalTok{(}
    \StringTok{"Term"}\NormalTok{, }\StringTok{"Estimate"}\NormalTok{,}
    \StringTok{"95\% CI Lower"}\NormalTok{, }\StringTok{"95\% CI Upper"}
\NormalTok{  ),}
  \AttributeTok{caption =} \StringTok{"Linear model coefficients ranked by}
\StringTok{    absolute effect size."}
\NormalTok{)}
\end{Highlighting}
\end{Shaded}

Both models agree that flipper length and species are the most important
determinants of body mass. The linear model additionally provides the
direction and magnitude of each effect (e.g., each millimetre of flipper
length adds approximately 15-20 grams), which the random forest
importance scores alone cannot convey.

\section{Model Selection Guidelines}\label{model-selection-guidelines}

The following framework maps common project requirements to a
recommended model class.

\begin{Shaded}
\begin{Highlighting}[]
\NormalTok{guidelines }\OtherTok{\textless{}{-}} \FunctionTok{tibble}\NormalTok{(}
  \AttributeTok{Requirement =} \FunctionTok{c}\NormalTok{(}
    \StringTok{"Regulatory or peer review"}\NormalTok{,}
    \StringTok{"Small sample (n \textless{} 500)"}\NormalTok{,}
    \StringTok{"Coefficient interpretation needed"}\NormalTok{,}
    \StringTok{"Maximum predictive accuracy"}\NormalTok{,}
    \StringTok{"Complex feature interactions"}\NormalTok{,}
    \StringTok{"Quick prototyping"}
\NormalTok{  ),}
  \AttributeTok{Recommendation =} \FunctionTok{c}\NormalTok{(}
    \StringTok{"Linear model"}\NormalTok{,}
    \StringTok{"Linear model"}\NormalTok{,}
    \StringTok{"Linear model"}\NormalTok{,}
    \StringTok{"Random forest"}\NormalTok{,}
    \StringTok{"Random forest"}\NormalTok{,}
    \StringTok{"Linear model"}
\NormalTok{  ),}
  \AttributeTok{Rationale =} \FunctionTok{c}\NormalTok{(}
    \StringTok{"Transparent, auditable logic"}\NormalTok{,}
    \StringTok{"Fewer parameters, lower variance"}\NormalTok{,}
    \StringTok{"Direct effect estimates with CIs"}\NormalTok{,}
    \StringTok{"Flexible, captures non{-}linearity"}\NormalTok{,}
    \StringTok{"Automatic interaction detection"}\NormalTok{,}
    \StringTok{"Fast fitting, minimal tuning"}
\NormalTok{  )}
\NormalTok{)}

\FunctionTok{kable}\NormalTok{(}
\NormalTok{  guidelines,}
  \AttributeTok{col.names =} \FunctionTok{c}\NormalTok{(}
    \StringTok{"Requirement"}\NormalTok{, }\StringTok{"Recommendation"}\NormalTok{,}
    \StringTok{"Rationale"}
\NormalTok{  ),}
  \AttributeTok{caption =} \StringTok{"Model selection guidelines based on}
\StringTok{    project requirements."}
\NormalTok{)}
\end{Highlighting}
\end{Shaded}

For the Palmer Penguins dataset, the linear model satisfies the first
three criteria and falls only marginally short on predictive accuracy.
The random forest is preferred only when the application demands maximum
predictive performance and the analyst is willing to forgo direct
coefficient interpretation.

\section{Practical Prediction
Comparison}\label{practical-prediction-comparison}

\begin{Shaded}
\begin{Highlighting}[]
\NormalTok{new\_penguin }\OtherTok{\textless{}{-}} \FunctionTok{data.frame}\NormalTok{(}
  \AttributeTok{species =} \StringTok{"Gentoo"}\NormalTok{,}
  \AttributeTok{flipper\_length\_mm =} \DecValTok{220}\NormalTok{,}
  \AttributeTok{bill\_length\_mm =} \DecValTok{47}\NormalTok{,}
  \AttributeTok{bill\_depth\_mm =} \DecValTok{15}\NormalTok{,}
  \AttributeTok{sex =} \StringTok{"male"}\NormalTok{,}
  \AttributeTok{island =} \StringTok{"Biscoe"}
\NormalTok{)}

\NormalTok{linear\_pred }\OtherTok{\textless{}{-}} \FunctionTok{predict}\NormalTok{(}
\NormalTok{  cv\_linear, }\AttributeTok{newdata =}\NormalTok{ new\_penguin}
\NormalTok{)}
\NormalTok{rf\_pred }\OtherTok{\textless{}{-}} \FunctionTok{predict}\NormalTok{(}
\NormalTok{  cv\_rf, }\AttributeTok{newdata =}\NormalTok{ new\_penguin}
\NormalTok{)}

\NormalTok{pred\_comparison }\OtherTok{\textless{}{-}} \FunctionTok{tibble}\NormalTok{(}
  \AttributeTok{Model =} \FunctionTok{c}\NormalTok{(}\StringTok{"Linear"}\NormalTok{, }\StringTok{"Random Forest"}\NormalTok{),}
  \AttributeTok{Predicted\_mass\_g =} \FunctionTok{c}\NormalTok{(linear\_pred, rf\_pred),}
  \AttributeTok{Difference\_g =} \FunctionTok{c}\NormalTok{(}
    \ConstantTok{NA}\NormalTok{, rf\_pred }\SpecialCharTok{{-}}\NormalTok{ linear\_pred}
\NormalTok{  )}
\NormalTok{)}

\FunctionTok{kable}\NormalTok{(}
\NormalTok{  pred\_comparison,}
  \AttributeTok{digits =} \DecValTok{0}\NormalTok{,}
  \AttributeTok{col.names =} \FunctionTok{c}\NormalTok{(}
    \StringTok{"Model"}\NormalTok{, }\StringTok{"Predicted Mass (g)"}\NormalTok{, }\StringTok{"Difference (g)"}
\NormalTok{  ),}
  \AttributeTok{caption =} \StringTok{"Predicted body mass for a hypothetical}
\StringTok{    male Gentoo penguin from Biscoe Island."}
\NormalTok{)}
\end{Highlighting}
\end{Shaded}

The two models produce predictions that differ by a modest amount, well
within the typical measurement uncertainty for field-weighed penguins.

\section{Series Summary}\label{series-summary}

This five-part series traced the full modelling workflow from
exploratory data analysis to model selection.

\begin{Shaded}
\begin{Highlighting}[]
\NormalTok{series\_summary }\OtherTok{\textless{}{-}} \FunctionTok{tibble}\NormalTok{(}
  \AttributeTok{Part =} \FunctionTok{paste}\NormalTok{(}\StringTok{"Part"}\NormalTok{, }\DecValTok{1}\SpecialCharTok{:}\DecValTok{5}\NormalTok{),}
  \AttributeTok{Topic =} \FunctionTok{c}\NormalTok{(}
    \StringTok{"EDA and Simple Regression"}\NormalTok{,}
    \StringTok{"Multiple Regression and Species"}\NormalTok{,}
    \StringTok{"Cross{-}Validation and Model Comparison"}\NormalTok{,}
    \StringTok{"Model Diagnostics and Interpretation"}\NormalTok{,}
    \StringTok{"Random Forest versus Linear Models"}
\NormalTok{  ),}
  \AttributeTok{Key\_result =} \FunctionTok{c}\NormalTok{(}
    \StringTok{"R{-}squared = 0.759 (flipper only)"}\NormalTok{,}
    \StringTok{"R{-}squared = 0.863 (added species)"}\NormalTok{,}
    \StringTok{"CV confirms robust generalisation"}\NormalTok{,}
    \StringTok{"All four assumptions satisfied"}\NormalTok{,}
    \StringTok{"RF gains only \textasciitilde{}2 pp over linear"}
\NormalTok{  )}
\NormalTok{)}

\FunctionTok{kable}\NormalTok{(}
\NormalTok{  series\_summary,}
  \AttributeTok{col.names =} \FunctionTok{c}\NormalTok{(}\StringTok{"Part"}\NormalTok{, }\StringTok{"Topic"}\NormalTok{, }\StringTok{"Key Result"}\NormalTok{),}
  \AttributeTok{caption =} \StringTok{"Summary of findings across all five}
\StringTok{    parts of the Palmer Penguins series."}
\NormalTok{)}
\end{Highlighting}
\end{Shaded}

The central narrative is one of diminishing returns from complexity. The
largest performance gain came from incorporating biological context
(species) in Part 2. Every subsequent addition (polynomial terms, random
forests, additional predictors) yielded only marginal improvements.

\subsection{Things to Watch Out For}\label{things-to-watch-out-for}

\begin{enumerate}
\def\labelenumi{\arabic{enumi}.}
\item
  \textbf{Unfair comparisons.} Giving one model access to additional
  predictors that the other lacks introduces a confound. We did this
  deliberately to give the random forest every advantage, but in a
  formal comparison both models should use the same predictor set.
\item
  \textbf{Overstating small differences.} A two- percentage-point
  improvement in R-squared may be statistically indistinguishable from
  zero given the sample size. Always check fold-by-fold variability
  before drawing conclusions.
\item
  \textbf{Feature importance is not causation.} The random forest's
  variable importance ranking reflects predictive utility, not causal
  effect. Flipper length predicts body mass well, but lengthening a
  penguin's flippers will not make it heavier.
\item
  \textbf{Reproducibility of random forests.} Different random seeds
  produce different forests and therefore slightly different importance
  rankings. Always set a seed and report it.
\item
  \textbf{Communication burden.} A linear model can be explained in a
  single equation; a random forest requires diagrams, importance plots,
  and partial dependence plots. Consider the audience before choosing
  the model.
\end{enumerate}

\begin{figure}[H]

{\centering \pandocbounded{\includegraphics[keepaspectratio,alt={Wide Antarctic panorama with penguins visible in the foreground and mountains behind}]{media/images/ambiance3.jpg}}

}

\caption{Antarctic landscape with penguins in the distance}

\end{figure}%

\emph{Closing visual before the summary sections.}

\subsection{Lessons Learnt}\label{lessons-learnt}

\subsubsection{Conceptual Understanding}\label{conceptual-understanding}

\begin{itemize}
\tightlist
\item
  The random forest outperformed the linear model by approximately two
  percentage points of R-squared, a difference that is modest relative
  to the loss of interpretability.
\item
  Both models identified flipper length and species as the dominant
  predictors, confirming the biological plausibility of the linear
  model's structure.
\item
  The largest single improvement in the entire series came from adding
  species to the model in Part 2 (R-squared from 0.76 to 0.86), not from
  adopting a more complex algorithm.
\item
  For well-behaved ecological data with moderate sample sizes, linear
  models with appropriate biological covariates can match or closely
  approximate ensemble methods.
\end{itemize}

\subsubsection{Technical Skills}\label{technical-skills}

\begin{itemize}
\tightlist
\item
  The \texttt{caret} package enables fair model comparison by enforcing
  identical cross-validation folds across different algorithms via
  \texttt{trainControl}.
\item
  \texttt{varImp()} provides a unified interface for extracting variable
  importance from both linear and ensemble models, facilitating direct
  comparison.
\item
  Setting \texttt{importance\ =\ TRUE} in \texttt{randomForest} computes
  permutation-based importance, which is more reliable than
  impurity-based measures for correlated predictors.
\item
  The \texttt{broom} package's \texttt{tidy()} and \texttt{glance()}
  functions make it straightforward to extract publication-ready
  coefficient tables from linear models.
\end{itemize}

\subsubsection{Gotchas and Pitfalls}\label{gotchas-and-pitfalls}

\begin{itemize}
\tightlist
\item
  The \texttt{caret} package's \texttt{train()} function defaults to
  tuning the \texttt{mtry} hyperparameter for random forests, which can
  produce confusing output if the analyst expects a single set of
  results.
\item
  Random forest predictions from \texttt{caret} objects require that the
  new data include all predictor columns, including factor levels
  present in the training data.
\item
  Large \texttt{ntree} values (e.g., 500 or 1000) improve stability but
  increase computation time; 500 trees is generally sufficient for
  datasets of this size.
\item
  The \texttt{caret} package loads many dependencies; in a reproducible
  pipeline, pin all package versions with \texttt{renv}.
\end{itemize}

\subsection{Limitations}\label{limitations}

\begin{itemize}
\tightlist
\item
  The random forest was given two additional predictors (sex and island)
  that the linear model did not use, making the comparison deliberately
  generous to the ensemble method.
\item
  The dataset contains only 333 complete observations from a single
  geographic location; results may not generalise to larger or more
  diverse penguin populations.
\item
  The temporal window (2007-2009) is narrow; both models assume
  stationary relationships that may shift under changing environmental
  conditions.
\item
  We did not tune random forest hyperparameters beyond the default
  \texttt{mtry} grid, which means the ensemble's performance ceiling may
  be slightly higher.
\item
  Partial dependence plots and interaction effects were not examined;
  doing so would provide a richer comparison of the two model classes.
\item
  The analysis does not account for potential measurement error in the
  morphometric variables.
\end{itemize}

\subsection{Opportunities for
Improvement}\label{opportunities-for-improvement}

\begin{enumerate}
\def\labelenumi{\arabic{enumi}.}
\tightlist
\item
  Conduct a fair comparison in which both models use the same predictor
  set (including sex and island in the linear model).
\item
  Tune random forest hyperparameters (mtry, ntree, nodesize) using
  nested cross-validation.
\item
  Add gradient-boosted trees (e.g., \texttt{xgboost}) as a third
  competitor to broaden the comparison.
\item
  Generate partial dependence plots for the random forest to visualise
  the shape of each variable's marginal effect.
\item
  Fit a mixed-effects model with species as a random effect to account
  for population-level variation more formally.
\item
  Extend the comparison to classification tasks (e.g., predicting
  species from morphometrics) to test whether the tradeoff differs for
  categorical outcomes.
\end{enumerate}

\section{Wrapping Up}\label{wrapping-up}

This series began with a simple scatter plot of flipper length versus
body mass and ended with a head-to-head comparison between a linear
model and a random forest. The central finding is that, for the Palmer
Penguins dataset, the species-aware linear model provides an excellent
balance of performance and interpretability. The random forest gains
only about two percentage points of R-squared while sacrificing the
transparency that makes linear models so valuable in scientific
communication.

The principal lesson from this exercise is that model selection is not
purely a technical decision. The right model depends on the audience,
the regulatory context, and the relative importance of explanation
versus prediction. For peer-reviewed ecological research, a
well-validated linear model with clear coefficient interpretations will
almost always be preferred over an opaque ensemble that offers marginal
predictive gains.

For anyone beginning a modelling journey, three pieces of advice are
offered. First, start simple and add complexity only when the data
demand it. Second, always cross-validate before claiming that a model
generalises. Third, run diagnostics to confirm that the model's
assumptions are satisfied; good predictions from a misspecified model
are a coincidence, not a guarantee.

Main takeaways:

\begin{itemize}
\tightlist
\item
  The random forest outperformed the linear model by only about two
  percentage points of R-squared.
\item
  Both models agree on the importance of flipper length and species.
\item
  The largest gain came from biological context (species), not
  algorithmic sophistication.
\item
  For ecological data of this type, the linear model is the recommended
  default.
\end{itemize}

\section{See Also}\label{see-also}

\textbf{Series posts:}

\begin{itemize}
\tightlist
\item
  \href{../08-palmerpenguinspart1/\%20palmerpenguinspart1/}{Part 1: EDA
  and Simple Regression}
\item
  \href{../09-palmerpenguinspart2/\%20palmerpenguinspart2/}{Part 2:
  Multiple Regression and Species Effects}
\item
  \href{../10-palmerpenguinspart3/\%20palmerpenguinspart3/}{Part 3:
  Cross-Validation and Model Comparison}
\item
  \href{../11-palmerpenguinspart4/\%20palmerpenguinspart4/}{Part 4:
  Model Diagnostics and Interpretation}
\end{itemize}

\textbf{Key resources:}

\begin{itemize}
\tightlist
\item
  Breiman, L. (2001). Random Forests. \emph{Machine Learning}, 45(1),
  5-32.
\item
  James, G., Witten, D., Hastie, T., \& Tibshirani, R. (2021). \emph{An
  Introduction to Statistical Learning} (2nd ed.). Springer.
  \url{https://www.statlearning.com}
\item
  Horst, A. M., Hill, A. P., \& Gorman, K. B. (2020). palmerpenguins:
  Palmer Archipelago (Antarctica) penguin data. R package.
  \url{https://allisonhorst.github.io/palmerpenguins/}
\item
  Kuhn, M. (2008). Building predictive models in R using the caret
  package. \emph{Journal of Statistical Software}, 28(5).
\item
  Molnar, C. (2022). \emph{Interpretable Machine Learning} (2nd ed.).
  \url{https://christophm.github.io/interpretable-ml-book/}
\end{itemize}

\section{Reproducibility}\label{reproducibility}

\textbf{Data source:} \texttt{palmerpenguins} R package (built-in
dataset, no external download required).

\textbf{Pipeline commands:}

\begin{Shaded}
\begin{Highlighting}[]
\ExtensionTok{Rscript}\NormalTok{ analysis/scripts/01\_prepare\_data.R}
\ExtensionTok{Rscript}\NormalTok{ analysis/scripts/02\_fit\_models.R}
\ExtensionTok{Rscript}\NormalTok{ analysis/scripts/03\_generate\_figures.R}
\ExtensionTok{quarto}\NormalTok{ render index.qmd}
\end{Highlighting}
\end{Shaded}

\textbf{Session information:}

\section{Feedback}\label{feedback}

Corrections, suggestions, and questions are welcome. Please open an
issue or pull request on the \href{https://github.com/mygit}{GitHub
repository} or send an email to user@example.com.




\end{document}
